% ---
% title: Lost cards: aces
% tags: generalized-bayes, total-probability, equally-likely
% ---

\begin{problem}
	Two cards have been lost from a deck of 52 cards, and we do not know which
	ones. We draw 4 of the remaining 50 cards at random.
	(a) Find the probability that there are exactly 2 aces among the 4 drawn
	cards.
	(b) There are exactly 2 aces among the 4 drawn cards. What is the
	probability that both lost cards were aces?
\end{problem}

\begin{solution}
	\idea The result of the draw depends on how many aces were lost. If we knew
	that number, the draw would be a simple equally-likely problem. So we split
	the possibilities by the number of lost aces, use the law of total
	probability in (a), and turn the conditioning around with the generalized
	Bayes formula in (b), as in Exercise~\exref{three-cooks}.

	\textbf{The partition.} Let $B_0$, $B_1$, $B_2$ be the events that $0$, $1$
	or $2$ aces were lost. Exactly one of them happens, so they are mutually
	exclusive and their union is $\Om$. The two lost cards are a random pair of
	the 52 cards, and every pair is equally likely. There are
	$\binom{52}{2}=1326$ pairs. A pair with no ace uses 2 of the 48 non-aces, a
	pair with one ace uses 1 of the 4 aces and 1 of the 48 non-aces, and a pair
	with two aces uses 2 of the 4 aces:
	\[ P(B_0)=\frac{\binom{48}{2}}{1326}=\frac{1128}{1326},\qquad
		P(B_1)=\frac{4\cdot48}{1326}=\frac{192}{1326},\qquad
		P(B_2)=\frac{\binom{4}{2}}{1326}=\frac{6}{1326}. \]
	Check: $1128+192+6=1326$.

	\textbf{The draw, for a known loss.} Let $A$ be the event that exactly 2 of
	the 4 drawn cards are aces. If $B_i$ happened, then $4-i$ aces and $46+i$
	non-aces are left among the 50 cards. Every set of 4 of these 50 cards is
	equally likely, and there are $\binom{50}{4}=230300$ such sets. For a
	favorable set we choose 2 of the remaining aces and 2 of the remaining
	non-aces:
	\[ P(A\mid B_0)=\frac{\binom{4}{2}\binom{46}{2}}{230300}
		=\frac{6\cdot1035}{230300}=\frac{6210}{230300}, \]
	\[ P(A\mid B_1)=\frac{\binom{3}{2}\binom{47}{2}}{230300}
		=\frac{3\cdot1081}{230300}=\frac{3243}{230300}, \]
	\[ P(A\mid B_2)=\frac{\binom{2}{2}\binom{48}{2}}{230300}
		=\frac{1\cdot1128}{230300}=\frac{1128}{230300}. \]
	The fewer aces were lost, the more aces are left, so the more likely it is
	to draw two of them.

	\step{(a)} The event $A$ is the union of the disjoint events $AB_0$,
	$AB_1$, $AB_2$, so their probabilities add up. By the multiplication rule,
	$P(AB_i)=P(B_i)P(A\mid B_i)$. This is the law of total probability:
	\[ P(A)=P(B_0)P(A\mid B_0)+P(B_1)P(A\mid B_1)+P(B_2)P(A\mid B_2). \]
	All three terms have the same denominator $1326\cdot230300=305\,377\,800$,
	so we add the numerators:
	\[ P(A)=\frac{1128\cdot6210+192\cdot3243+6\cdot1128}{305\,377\,800}
		=\frac{7\,004\,880+622\,656+6\,768}{305\,377\,800}
		=\frac{7\,634\,304}{305\,377\,800}. \]
	This fraction simplifies to
	\[ P(A)=\frac{6768}{270725}\approx0.0250. \]

	\step{(b)} Now we know that $A$ happened, and we ask for $P(B_2\mid A)$. By
	the definition of conditional probability, the numerator is the third term
	of the sum in (a), and the denominator is the whole sum. This is the
	generalized Bayes formula:
	\[ P(B_2\mid A)=\frac{P(B_2)P(A\mid B_2)}{P(A)}
		=\frac{6\cdot1128}{7\,634\,304}=\frac{6768}{7\,634\,304}
		=\frac{1}{1128}\approx0.00089. \]
	(The common denominator $305\,377\,800$ cancels.) For comparison, before
	the draw $P(B_2)=\frac{6}{1326}=\frac{1}{221}\approx0.0045$. Seeing two aces
	among the drawn cards makes it less likely that two aces were lost: if two
	aces were missing, only 2 would be left in the deck, and drawing both of
	them is unlikely.

	\begin{answer}
		(a) $P(A)=\frac{6768}{270725}\approx0.0250$ \quad
		(b) $P(B_2\mid A)=\frac{1}{1128}\approx0.00089$
	\end{answer}
	\begin{remark}
		The answer to (a) is exactly the probability of drawing 2 aces among 4
		cards from a full deck:
		\[ \frac{\binom{4}{2}\binom{48}{2}}{\binom{52}{4}}
			=\frac{6\cdot1128}{270725}=\frac{6768}{270725}. \]
		This is no accident. The lost cards are just two more cards taken at
		random, and nobody looks at them. Taking out 2 unseen cards and then
		drawing 4 is the same as drawing 4 cards from the full deck: every set of
		4 cards has the same chance either way. Losing cards changes nothing as
		long as we do not know which cards were lost. The information arrives
		only in (b), where we use the observed draw to learn about the lost
		cards.
		\textbf{Common mistake:} answering (b) with $P(B_2)=\frac{1}{221}$. That
		is the probability before the draw; the question asks for the
		probability after we have seen two aces.
	\end{remark}
\end{solution}
